\section{Week 15: Fixed All-Six Actual Encoder Success}

Week~15 closes \claimid{OBL-RANS-ACTUAL-SUCCESS-WITNESS} as
\status{proved for the fixed all-six input, Hoare partial correctness}.  The
result excludes the actual encoder's failure return for every terminating run
from the fixed precondition.  It does not establish termination or a
probability-one reachability theorem.

\subsection*{Zero HBZ to the actual all-six prepare result}

\theoryname{Mode2RansAllSixBudget} defines \procname{zero_hbz} by initializing
every byte of the 8192-byte coefficient array to zero.  The lemma
\procname{zero_hbz_get32} opens the four-byte \texttt{BArray8192.get32} load
and proves every coefficient word is \texttt{W32.zero}; canonicality is then
derived by \procname{zero_hbz_canonical}.  Thus canonicality is not an input
assumption of the final witness theorem.

The actual focused \procname{_encode_hb_z1_prepare} theorem yields
\texttt{bad=0}, the prepared-prefix relation, and its byte-tail frame.
\procname{prepared_zero_hbz_is_all_six} unfolds each prepared coefficient,
uses the zero word load, and proves that all first 1024 symbols equal 6.
Consequently \procname{actual_focused_encode_hb_z1_prepare_zero_hbz} connects
the concrete zero coefficient input to the all-six predicate without an
arbitrary prepared-symbol witness.

\subsection*{A coarse capacity proof}

The literal table certificate gives
\[
 f(6)=398,
 \qquad x_{\max}(6)=2{,}097{,}152\cdot398=834{,}666{,}496.
\]
For every state in the actual encoder range
$2^{23}\leq x<2^{31}$, integer division gives
$x\mathbin{\mathrm{div}}256<x_{\max}(6)$.  The lemma
\procname{symbol6_normalization_len_le1} therefore proves
$0\leq\ell_6(x)\leq1$: symbol 6 emits at most one normalization byte.

The first four all-six suffix states are checked inside EasyCrypt by unfolding
the certified fast/math update:
\[
 2^{23},\ 21{,}582{,}496,\ 55{,}528{,}910,\
 142{,}868{,}116,\ 367{,}580{,}518.
\]
The first four inputs are below $x_{\max}(6)$ before their updates, so
\procname{all_six_first_four_no_normalization} proves zero emitted bytes for
each of these four prefixes.  Induction then combines four zero-emission
steps with at most one byte for every remaining step:
\[
 \left|\mathsf{encodeTrace}(6^n).\mathsf{bytes}\right|
 \leq \max(0,n-4).
\]
At $n=1024$, \procname{all_six_trace_fits_mode2} derives at most 1020
normalization bytes and a total serialized trace size between 4 and 1024.
No externally computed exact trace length is used as proof evidence.

\subsection*{Failure guard and actual success}

The strengthened Week~11 closure retains the actual failure cause rather than
only \texttt{bad}\,$\ne0$.  The theorem
\procname{actual_rans_encode_failure_trace_cause} proves that a nonzero bad
return requires more than 1020 normalization bytes.  This certificate is
threaded from the generated \texttt{off < 4} guard through the existing
outer/inner tail invariants; the encoder loops are not copied or re-proved in
a new theory.

\procname{actual_rans_encode_all_six_success} targets
\procname{RansEncodeTarget.M._rans_encode} directly.  Its precondition binds
the actual arrays, literal mode-2 table, count 1024, and the all-six prefix; it
does not assume success, buffer fit, trace size, or a returned offset.  The
pure 1020-byte upper bound contradicts the actual failure certificate.  The
remaining branch has \texttt{bad=0}, offset in $[0,1020]$, serialized size in
$[4,1024]$, the exact all-six trace segment, and the initial prefix frame.

\subsection*{Full and production boundaries}

\breakablett{actual_encode_hb_z1_full_zero_success} opens the actual
\breakablett{HbzFullEncodeTarget.M._encode_hb_z1_full}.  It derives zero-input
canonicality, prepare success, and all-six symbols internally; calls the
actual rANS encoder theorem; and follows the actual suffix-copy branch.  Its
postcondition gives a nonzero returned size, the 4--1024 bound, exact trace,
output-tail frame, and the actual prepare-call symbol witness.

The exact focused/production procedure equivalence yields
\breakablett{signature_pack_hbz_zero_success_mode2} for
\breakablett{SignaturePackMode2Target.M._encode_hb_z1_full}.  The additional
corollary \breakablett{signature_pack_unpack_hbz_zero_success_mode2} eliminates
the failure disjunct of the existing production pack/unpack harness using the
new encoder theorem, then obtains decoder execution, decoder \texttt{bad=0},
the zero coefficient prefix, and both coefficient and encoded tail frames.

\subsection*{Boundary and next target}

All Week~15 claims remain unary Hoare partial correctness.  No
\texttt{phoare}=1 or losslessness theorem was compiled, so fixed-input
termination and actual non-vacuous reachability remain \status{partial}.
General encoder termination, all-canonical-input success, malformed-input
rejection, encoding zero-loss, and security are not claimed.  The Week~16
single recommended target is \claimid{OBL-SIG-H-ENCODE-DECODE}: the actual
mode-2 \(h\)-suffix codec leaf/table/loop/full-wrapper inverse.
